Three dimensional (3D) data have become a key modality for representing and interacting with the physical world. Modern vehicles and roadside infrastructure are equipped with 3D sensors that capture the geometry of traffic participants and surroundings. Meanwhile, the growing deployment of vehicular communication technologies enables these vehicles and infrastructure elements to exchange 3D data with one another and with cloud servers in real time. Together, advances in 3D sensing and vehicular connectivity are transforming how spatial information is captured, shared, and utilized in vehicular systems.

This dissertation explores how 3D data can be elevated to a shared and networked resource that is collaboratively sensed, processed, and delivered across vehicles, infrastructure, and cloud systems. It presents three systems spanning the full networked 3D data pipeline. \recap addresses the \textit{reconstruction} phase by enabling vehicles to offload 3D sensor data to the cloud for high-fidelity multi-view fusion under dynamic and sparsely overlapping observations. \cip targets the \textit{interpretation} phase, demonstrating real-time cooperative infrastructure perception by fusing data from multiple roadside LiDARs to track and estimate the motion of traffic participants. \mars extends the pipeline to the \textit{delivery} phase by streaming immersive 3D video content from the cloud to in-vehicle augmented reality (AR) displays, adapting to vehicle motion and dynamic network conditions to provide spatially coherent AR experiences.

Together, these contributions demonstrate that networked 3D data can be effectively sensed, reconstructed, interpreted, and delivered in connected vehicular environments, enhancing perception, improving traffic safety and efficiency, and enabling next-generation in-vehicle experiences.
